%==============================================================================
%  figures/solver_numerics.tex
%  Structural differences between the three solvers: geometry representation
%  and discretization. These differences are the source of the model
%  discrepancy exploited in the cross-model study. Adapted from presentation.tex.
%  \input in the methodology forward-models subsection (see tab:solver-differences).
%==============================================================================
\begin{figure}[tb]
  \centering
  % ---- (a) uDALES: staggered C-grid + cut cells ----------------------------
  \begin{subfigure}[b]{\linewidth}
    \centering
    \resizebox{\linewidth}{!}{%
    \begin{tikzpicture}[font=\scriptsize,>=latex]
      % LEFT: staggered Arakawa C-grid
      \draw[uaCharcoal!45,line width=0.6pt,step=1.8] (0,0) grid (3.6,3.6);
      \draw[uaCharcoal!70,line width=0.9pt] (0,0) rectangle (3.6,3.6);
      \foreach \x in {0.9,2.7}\foreach \y in {0.9,2.7}
        \fill[uaCharcoal] (\x,\y) circle (2.6pt);
      \node[text=uaCharcoal,anchor=west] at (0.98,1.12) {$p$};
      \foreach \x in {0,1.8,3.6}\foreach \y in {0.9,2.7}
        \draw[->,uaOrange,line width=1.1pt] (\x-0.28,\y) -- (\x+0.28,\y);
      \node[text=uaOrange] at (1.8,3.0) {$u$};
      \foreach \y in {0,1.8,3.6}\foreach \x in {0.9,2.7}
        \draw[->,uaTeal,line width=1.1pt] (\x,\y-0.28) -- (\x,\y+0.28);
      \node[text=uaTeal] at (0.62,1.95) {$w$};
      \node[font=\footnotesize\bfseries,text=uaCharcoal] at (1.8,4.05)
        {Staggered Arakawa C-grid};
      % RIGHT: cut cells
      \begin{scope}[xshift=7.0cm]
        \fill[uaCharcoal!22] (1.35,3.6) -- (4.5,0.45) -- (4.5,3.6) -- cycle;
        \draw[uaCharcoal!45,line width=0.5pt,step=0.9] (0,0) grid (4.5,3.6);
        \draw[uaCharcoal!70,line width=0.9pt] (0,0) rectangle (4.5,3.6);
        \foreach \c in {(0.9,2.7),(1.8,2.7),(1.8,1.8),(2.7,1.8),(2.7,0.9),(3.6,0.9),(3.6,0)}
          \draw[uaOrange,line width=1.1pt] \c rectangle ++(0.9,0.9);
        \draw[uaCharcoal,line width=1.6pt] (1.35,3.6) -- (4.5,0.45);
        \node[text=uaCharcoal,rotate=-45] at (3.45,2.0) {\bfseries facade};
        \node[font=\footnotesize\bfseries,text=uaOrange] at (2.25,4.05)
          {Conservative cut cells};
      \end{scope}
    \end{tikzpicture}}
    \caption{\pyudales: staggered C-grid with conservative cut-cell walls.}
    \label{fig:solver:udales}
  \end{subfigure}

  \vspace{6pt}
  % ---- (b) PALM: staircase walls + upwind advection ------------------------
  \begin{subfigure}[b]{\linewidth}
    \centering
    \resizebox{\linewidth}{!}{%
    \begin{tikzpicture}[font=\scriptsize,>=latex]
      % LEFT: staircase topography
      \fill[uaCharcoal!60] (0,0) rectangle (0.9,0.9);
      \fill[uaCharcoal!60] (0.9,0) rectangle (1.8,1.8);
      \fill[uaCharcoal!60] (1.8,0) rectangle (2.7,1.8);
      \fill[uaCharcoal!60] (2.7,0) rectangle (3.6,2.7);
      \fill[uaCharcoal!60] (3.6,0) rectangle (4.5,2.7);
      \fill[uaAmber!45] (0,0.9) rectangle (0.9,1.8);
      \fill[uaAmber!45] (0.9,1.8) rectangle (2.7,2.7);
      \fill[uaAmber!45] (2.7,2.7) rectangle (4.5,3.6);
      \draw[uaCharcoal!45,line width=0.5pt,step=0.9] (0,0) grid (4.5,3.6);
      \draw[uaCharcoal!70,line width=0.9pt] (0,0) rectangle (4.5,3.6);
      \draw[uaOrange,line width=1.4pt,densely dashed] (0,1.0) -- (4.5,3.25);
      \node[text=uaOrange,rotate=26] at (1.6,2.05) {\bfseries true surface};
      \node[font=\footnotesize\bfseries,text=uaCharcoal] at (2.25,4.05)
        {Mask method: staircase walls};
      % RIGHT: 5th-order upwind stencil
      \begin{scope}[xshift=7.2cm]
        \draw[->,uaTeal,line width=1.4pt] (0.4,3.3) -- (2.2,3.3)
          node[right,text=uaTeal]{\bfseries wind};
        \draw[uaCharcoal!45,line width=0.6pt] (0,1.2) grid[step=0.9] (5.4,2.1);
        \draw[uaCharcoal!70,line width=0.9pt] (0,1.2) rectangle (5.4,2.1);
        \foreach \x in {0.45,1.35,2.25,3.15}
          \fill[uaOrange] (\x,1.65) circle (2.6pt);
        \foreach \x in {4.05,4.95}
          \fill[uaOrange!40] (\x,1.65) circle (2.6pt);
        \draw[uaTeal,line width=2.0pt] (3.6,1.08) -- (3.6,2.22);
        \node[text=uaTeal,anchor=north] at (3.6,1.0) {\bfseries flux face};
        \draw[uaOrange,line width=0.9pt] (0.45,2.45) -- (0.45,2.6) -- (4.95,2.6)
          -- (4.95,2.45);
        \node[text=uaOrange] at (2.7,2.85) {6-point stencil, upwind-biased};
        \node[font=\footnotesize\bfseries,text=uaCharcoal] at (2.7,4.05)
          {5th-order upwind advection};
      \end{scope}
    \end{tikzpicture}}
    \caption{\pypalm: staircase (mask) topography with high-order upwind advection.}
    \label{fig:solver:palm}
  \end{subfigure}

  \vspace{6pt}
  % ---- (c) LBM: collide & stream on a uniform lattice ----------------------
  \begin{subfigure}[b]{\linewidth}
    \centering
    \resizebox{0.9\linewidth}{!}{%
    \begin{tikzpicture}[font=\scriptsize,>=latex]
      % LEFT: discrete velocity set (D2Q9 sketch)
      \draw[uaCharcoal!25,line width=0.5pt,step=1.1] (-1.1,-1.1) grid (1.1,1.1);
      \fill[uaCharcoal] (0,0) circle (3pt);
      \foreach \a in {0,90,180,270}
        \draw[->,uaTeal,line width=1.2pt] (0,0) -- (\a:1.0);
      \foreach \a in {45,135,225,315}
        \draw[->,uaTeal!60,line width=1.2pt] (0,0) -- (\a:1.35);
      \node[text=uaTeal] at (1.45,0.42) {$f_i$};
      \node[font=\footnotesize\bfseries,text=uaCharcoal] at (0,1.85)
        {Populations $f_i$};
      \node[text=uaGrey,align=center,text width=3.4cm] at (0,-1.95)
        {discrete velocity set\\ (D2Q9 sketch)};
      % MIDDLE: collide & stream
      \begin{scope}[xshift=4.6cm]
        \draw[uaCharcoal!25,line width=0.5pt,step=1.1] (-1.1,-1.1) grid (2.2,1.1);
        \foreach \x in {-1.1,0,1.1,2.2}\foreach \y in {-1.1,0,1.1}
          \fill[uaCharcoal!55] (\x,\y) circle (1.8pt);
        \fill[uaCharcoal] (0,0) circle (3pt);
        \foreach \t in {(1.1,0),(0,1.1),(-1.1,0),(0,-1.1),(1.1,1.1),(-1.1,1.1),(1.1,-1.1),(-1.1,-1.1)}
          \draw[->,uaOrange,line width=1.1pt,shorten >=3pt,shorten <=4pt] (0,0) -- \t;
        \node[font=\footnotesize\bfseries,text=uaOrange] at (0.55,1.85)
          {Collide \& stream};
        \node[text=uaGrey,align=center,text width=4.2cm] at (0.55,-1.95)
          {local relaxation, then hop\\ to neighbours -- no global solve};
      \end{scope}
      % RIGHT: bounce-back wall
      \begin{scope}[xshift=10.1cm]
        \fill[uaCharcoal!60] (-1.4,-1.5) rectangle (1.4,-1.1);
        \draw[uaCharcoal,line width=1.2pt] (-1.4,-1.1) -- (1.4,-1.1);
        \draw[uaCharcoal!25,line width=0.5pt,step=1.1] (-1.1,-1.1) grid (1.1,1.1);
        \fill[uaCharcoal] (0,0) circle (3pt);
        \draw[->,uaTeal,line width=1.2pt] (0.10,0.02) -- (-0.40,-0.98);
        \draw[->,uaOrange,line width=1.2pt,densely dashed] (-0.62,-0.98) -- (-0.12,0.02);
        \node[font=\footnotesize\bfseries,text=uaCharcoal] at (0,1.85)
          {Bounce-back walls};
        \node[text=uaGrey,align=center,text width=3.8cm] at (0,-2.35)
          {populations reflect\\ at solid nodes};
      \end{scope}
    \end{tikzpicture}}
    \caption{\lbmmodel: collide-and-stream on a uniform lattice with bounce-back
      walls (2D D2Q9 sketch; the solver uses D3Q27).}
    \label{fig:solver:lbm}
  \end{subfigure}
  \caption{The three solvers share the governing equations but differ in
    discretization, geometry representation, and near-wall treatment ---
    \pyudales\ (a) cut cells, \pypalm\ (b) staircase masks, \lbmmodel\ (c)
    bounce-back on a uniform lattice. These structural differences are the
    source of the \discrepancy\ probed by the \crossmodel\ experiments; see
    \tabref{tab:solver-differences}.}
  \label{fig:solver_numerics}
\end{figure}
